\documentclass[10pt,a4paper]{amsart}
\usepackage[margin=19mm]{geometry}
\usepackage{amsmath,amssymb,mathtools}
\usepackage[scr=rsfs,cal=euler]{mathalfa}
\usepackage{xfrac}
\usepackage[unicode,hidelinks,bookmarks=false]{hyperref}
\newcommand{\ie}{i.e.\ }
\newcommand{\cf}{cf.\ }
\newcommand{\resp}{respectively\ }
\newcommand{\Subsection}{\subsection}
\providecommand{\nobreakdash}{\nobreak\hbox{-}}
\setlength{\emergencystretch}{2em}
\multlinegap0pt
\usepackage{amssymb}
\usepackage{mathtools}
\newtheorem{hyp}{Hypothesis}
\newtheorem{proposition}{Proposition}
\newtheorem{pf}{Pf}
\def\R{\mathbb{R}}
\def\geq{\geqslant}
\def\leq{\leqslant}
\title{Convergence Rate of Euler--Maruyama Scheme to the Invariant Probability Measure Under Total Variation Distance for the SDEs}
\author{Yuke Wang, Yinna Ye}
\date{Source: arXiv:2505.04218v3 (CC BY 4.0)}
\begin{document}
\maketitle

\noindent\emph{Source and licence.} Convergence Rate of Euler--Maruyama Scheme to the Invariant Probability Measure Under Total Variation Distance for the SDEs. Version arXiv:2505.04218v3 (CC BY 4.0), distributed by arXiv under the Creative Commons Attribution 4.0 International licence, http://creativecommons.org/licenses/by/4.0/, as recorded in the arXiv licence metadata for this exact version and shown on the abstract page https://arxiv.org/abs/2505.04218v3. Full licence notice: \texttt{licenses/arxiv-2505.04218-CC-BY-4.0.txt}.\par\smallskip

\noindent\textit{Editorial scope.} Counted unit: the article's proposition aperiodic (source0.tex lines 251--261). The article's complete original proof of it is delivered with it (source0.tex lines 262--281, one \texttt{proof} environment of 20 lines). No mathematics is written, repaired or re-derived: the statement and the proof are the article's own lines, in the article's own notation, and no retained line is substituted or re-flowed.  Retained uncounted as the source-local context the proof needs: lines 246--248.  Nothing was omitted from any retained span, so the package carries no omission record.  No retained line contains a citation, so the wrapper carries no bibliography and no bibliography entry.  No retained line contains a figure environment, an image-inclusion command or drawing code, so no image is delivered and the package declares no asset dependency.  Editorial formatting only, in addition: an AMS \texttt{amsart} preamble that restates the article's own declarations and macros used by the retained lines, namely the environments \texttt{hyp}, \texttt{proposition}, \texttt{pf} on separate counters, together with the standard packages those lines need.  The retained environments print in this document as Hyp 1, Proposition 1, Pf 1 in order of appearance, whereas the article prints the same items with the numbering of its own full document.  The complete native source of the article as distributed in the arXiv e-print for this exact version is bundled unchanged as \texttt{sources/177728/source0.tex}.
\begin{hyp}\label{hyp2}\label{assump2}
	For any $x\in\R$ and $\eta\in(0,1)$, $|x+\eta g(x)|\leq c.$
\end{hyp}

\begin{proposition}\label{aperiodic}
	\begin{enumerate}
		\item[(1)] If $C$ is a compact subset of $\mathbb{R}$ such that $\mbox{Leb}(C)>0$,  then $C$ is an accessible $(1,\epsilon\nu)$-small set, where
		\begin{align}\label{Def-eps}
			\epsilon:=\frac{\mbox{Leb}(C)}{\sqrt{\eta}\,\sigma}\inf_{(x,y)\in C^2}\phi\left(\frac{y-x-\eta g(x)}{\sqrt{\eta}\sigma}\right),
		\end{align}
		$\nu(\cdot):=\mbox{Leb}(\cdot\,\cap C)/\mbox{Leb}(C)$, and~$\phi(x):=\frac{1}{\sqrt{2\pi}}\exp\{-x^2/2\}$ is the probability density function of standard normal distribution.
		\item[(2)] For any $\eta\in(0,1)$, the~Markov kernel $P_{\eta}$ is irreducible and strongly aperiodic.
		\item[(3)] Under Hypothesis \ref{hyp2}, the~state space $\mathbb{R}$ of the Markov chain $(\theta_k)_{k\geq0}$ is an $1$-small set.
	\end{enumerate}
\end{proposition}

\begin{pf}
	\begin{enumerate}
	\item[(1)] 	Suppose that $C$ is a compact subset of $\mathbb{R}$ such that $\mbox{Leb}(C)>0$. Recall that $P_{\eta}$ admits a Markov kernel $p_{\eta}(x,y)$, given~by
	$$p_{\eta}(x,y)=\frac{1}{\sqrt{\eta}\,\sigma}\phi\left(\frac{y-x-\eta g(x)}{\sqrt{\eta}\sigma} \right).$$
	Then for all $x\in C$ and $A\in\mathcal{B}(\R)$,		
	\begin{align*}
			P_{\eta}(x,A)=\int_A p_{\eta}(x,y) dy\geq\int_{A\cap C}p_{\eta}(x,y) dy\geq\,\epsilon\, \frac{\mbox{Leb}(A\cap C)}{\mbox{Leb}(C)},
	\end{align*}
	where $\epsilon\in(0,\,1]$ is defined in (\ref{Def-eps}). Evidently, $C$ is an $(1,\epsilon\nu)$-small set, with $\nu(\cdot)=\mbox{Leb}(\cdot\,\cap C)/\mbox{Leb}(C)$. Furthermore, $C$ is accessible since for all $x\in\R$,
	$$P_{\eta}(x,C)=\frac{1}{\sqrt{\eta}\,\sigma}\int_C\phi\left(\frac{y-x-\eta g(x)}{\sqrt{\eta}\sigma}\right)dy>0.$$
	\item[(2)] We have proved in (1) that any compact subset $C$ of $\R$ satisfying $\mbox{Leb}(C)>0$ is an accessible $(1,\mu)$-small set with $\mu(C)=\epsilon\nu(C)>0$. Hence, $P_{\eta}$ is irreducible and strongly aperiodic.
	\item[(3)] Assume Hypothesis \ref{hyp2}. Then we have that for a given non-empty set $A\in\mathcal{B}(\R)$, both $\inf_{x\in\R}(x+\eta g(x))$ and $\sup_{x\in\R}(x+\eta g(x))$ exist in $\R$; for a given $y\in A$,	
		$$
		\inf_{x\in\R}p_{\eta}(x,y)\geq\frac{1}{\sqrt{2\pi\eta}\sigma}\exp\left\{-\frac{1}{2\eta\sigma^2}\left[\left(y-\inf_{x\in\R}(x+\eta g(x))\right)^2\vee\left(y-\sup_{x\in\R}(x+\eta g(x))\right)^2\right]\right\}>0.
		$$
      Let $f(y)$ denote the function in the middle of the two inequalities above. Then for all $x\in\R$ and $A\in\mathcal{B}(\R)$,
      $$P_{\eta}(x,\,A)=\int_A p_{\eta}(x,y) dy\geq\mu(A),$$
      where $\mu(A)=\int_A f(y) dy$ is a Borel measure defined on $\mathcal{B}(\R)$. Therefore, the~state space $\R$ is an $(1,\mu)$-small set. \hfill$\square$
\end{enumerate}
\end{pf}

\end{document}
